\section{Сеточное решение уравнения Пуассона}
\subsection{Постановка задачи}

Будем рассматривать многомерное дифференциальное уравнение в области $\Omega$:
\begin{equation}
    \label{eq:poissonnd_lam}
    -\nabla\cdot\left(\lambda(\vec x)\nabla u\right) = f(\vec x), \quad \vec x \in \Omega
\end{equation}

Оператор в левой части (оператор Лапласа)
описывает физический процесс диффузии c коэффициентом диффузии $\lambda$.
Это уравнение (с нулевой правой частью) используют
в частности для расчёта распределения температуры в однородном твёрдом теле.
В этом случае коэффиент $\lambda$ называют коэффициентом теплопроводности,
а за счёт ненулевой $f$ можно задавать дополнительные
внутренние источники тепла.

В простом случае постоянного коэффициент диффузии ($\lambda = \const$) его
можно вынести из под дивергеции и отнести в правую часть. Тогда уравнение упростится до однородного вида:
\begin{equation}
    \label{eq:poissonnd}
    -\nabla^2 u = f(\vec x), \quad \vec x \in \Omega
\end{equation}

Далее на границе области расчёта $\partial \Omega$
рассмотрим несколько типов граничных условий.

\subsubsection{Граничные условия первого рода}
Также известны как граничные условия Дирихле.
На границе $\partial \Omega_{I}$ задано точное значение искомой функции $u^\Gamma$:
\begin{equation}
\label{eq:poissonnd_bc}
u = u^\Gamma(\vec x), \quad \vec x \in \partial\Omega_I
\end{equation}

В аналогии задачи теплопроводности это условие
можно трактовать как условие заданной на стенке температуры.

В случае, если на всей границе области поставлены условия второго рода: $\Omega = \Omega_{II}$,
задача Пуассона~\cref{eq:poissonnd_lam} не будет иметь единственного решения (решение будет определено
с точностью до константы). Для получения единственного решения следует
однозначно определить решение в произвольной точке области
$$
u(\vec x_0) = u^0.
$$

\subsubsection{Граничные условия второго рода}
Также известны как граничные условия Неймана.
На границе $\partial \Omega_{II}$ задано значение нормальной производной искомой функции $q$:
\begin{equation}
\label{eq:poissonnd_bc2}
-\lambda(\vec x)\dfr{u}{n} = q(\vec x), \quad \vec x \in \partial\Omega_{II}
\end{equation}

В аналогии задачи теплопроводности это условие можно трактовать как условие заданного на стенке теплового потока.

\subsubsection{Граничные условия третьего рода}
Также известны как граничные условия Робэна.
На границе $\partial \Omega_{III}$ задано линейное соотношение значений функции и нормальной производной:
\begin{equation}
\label{eq:poissonnd_bc3}
-\lambda(\vec x)\dfr{u}{n} = \alpha(\vec x) u + \beta(\vec x), \quad \vec x \in \partial\Omega_{III}.
\end{equation}

При постановке задачи теплопроводности это условие часто записывают в виде
$$
-\lambda(\vec x)\dfr{u}{n} = \alpha(\vec x)\left(u - u^0\right).
$$
где известное значение $u^0$ называют температурой окружающей стреды.
Такое условие называют условием конвективной теплопроводности или условием Ньютона--Рихмана.
Для приведения этого условия к исходному виду \cref{eq:poissonnd_bc3} достаточно положить $\beta = -\alpha u^0$.

\subsubsubsection{Об универсальности условий третьего рода}
Условия \cref{eq:poissonnd_bc,eq:poissonnd_bc2}
можно свести к условиям \cref{eq:poissonnd_bc3}
при правильном подборе коэффициентов $\alpha$ и $\beta$.
Так для условий второго рода нужно положить $\alpha=0$, $\beta=q$.
А для условий первого рода: $\alpha=\eps^{-1}$, $\beta = -\alpha u^\Gamma$,
где $\eps\to 0$ -- некоторое очень малое положительное число.

\subsubsection{Периодические граничные условия}
Необходимость в таких условиях
возникает при расчёте физических процессов
около периодических структур: решёток, лопастей, рядов скважин, оребрения нагревателя и т.п.
В этом случае из исходную большую область расчёта
представляют как бесконечную последовательность однотипных ячеек периодичности, 
в каждой из которых решения полностью идентичны (в более сложных вариантах -- сдвинуты на константу).

\begin{figure}[h!]
\centering
\includegraphics[width=0.4\linewidth]{periodic_bc.pdf}
\caption{Ячейка периодичности в задаче обтекания бесконечной решётки}
\label{fig:periodic_bc}
\end{figure}
На \figref{fig:periodic_bc} представлен пример области с выделенной ячейкой периодичности $\overline\Omega$ (незатенённая область).
Изолинии можно трактовать как изотермы решения задачи о нестационарном обтекании решётки нагревателя
(поле температур в этом случае описывается более сложным уравнением, чем \cref{eq:poissonnd_lam}
и приведено тут только для иллюстрации периодичности).

Пара периодических границ обозначена через $\partial\overline\Omega_P$ и $\partial\overline\Omega_P'$.
Пусть эти границы топологически экваивалентны, то есть для любой точки $\vec x\in\partial\overline\Omega_P$
cуществует взаимноодносзначная точка $\vec x'\in\partial\overline\Omega_P'$.
Для того, чтобы решение за этими границами точно соответствовало решению внутри ячейки периодичности необходимо
задать равенство значений и производных любого порядка:
\begin{equation}
\label{eq:poissonnd_bcp}
\left\{
\begin{array}{l}
    u(\vec x) = u(\vec x'), \\ [10pt]
    \left.\displaystyle\frac{\partial^k u}{\partial n^k}\right|_{\vec x} = -\left.\displaystyle\frac{\partial^k u}{\partial n^k}\right|_{\vec x'}
\end{array}
\right.
\vec x\in\partial\overline\Omega_P, \quad \vec x'\in\partial\overline\Omega_P', \quad \forall k.
\end{equation}
Здесь под $n$ подразумевается внешняя к ячейке периодичности нормаль, поэтому в правой части
условия для производных стоит минус.


\subsection{Метод конечных разностей}
\label{sec:poisson_fdm}
Будем рассмотривать задачу \cref{eq:poissonnd,eq:poissonnd_bc} в упрощённой одномерной постановке:
\begin{equation}
    \label{eq:poisson1d}
    -\ddfrq{u}{x} = f(x)
\end{equation}
в области $x\in[a,b]$.

Для аппроксимации производных методом конечных разностей будем использовать разложение Тейлора:
\begin{equation}
    \label{eq:taylor_series}
    u(x+h) = u(x) + h u'(x) + \frac{h^2}{2} u''(x) + \dots + \frac{h^{k}}{k!} u^{(k)}(x)
\end{equation}
из которой можно получить приближение первого (направленная разность) и второго (симметричная разность) порядков:
\begin{align}
    \label{eq:forwdiff}
    &u'(x) = \frac{u(x+h) - u(x)}{h} \underbrace{- \frac{h}{2}u''(x) - \frac{h^2}{6}u'''(x) - \dots}_{O(h)}, \\
    \label{eq:symdiff}
    &u'(x) = \frac{u(x+h/2) - u(x-h/2)}{h} \underbrace{-\frac{h^2}{24} u'''(x) - \frac{h^4}{1920} u^{(5)}(x) - \dots}_{O(h^2)}.
\end{align}

\subsubsection{Аппроксимация во внутренней области}

В области решения $[a,b]$ введём равномерную сетку из $N$ ячеек.
Шаг сетки будет равен $h=(b-a)/N$.
Узлы сетки запишем в виде сеточного вектора $\{x_i\}$ длины $N+1$, где $i=\overline{0,N}$.
Определим сеточный вектор $\{u_i\}$ неизвестных, элементы которого определяют значение искомого численного решения в $i$-ом узле сетки. 

Для аппроксимации оператора Лапласа из левой части исходного уравнения~\cref{eq:poisson1d} сначала
воспользуемся симметричной разностью~\cref{eq:symdiff} для снятия первой производной:
\begin{equation}
\label{eq:fdm_laplace_init}
u''(x) = \frac{u'(x+h/2) - u'(x-h/2)}{h} + O(h^2).
\end{equation}
Далее распишем первые производную с помощью той же формулы~\cref{eq:symdiff}:
\begin{align*}
& u'(x+h/2) = \frac{u(x+h) - u(x)}{h} - \frac{h^2}{24} u'''(x+h/2) + O(h^4)\\
& u'(x-h/2) = \frac{u(x) - u(x-h)}{h} - \frac{h^2}{24} u'''(x+h/2) + O(h^4).
\end{align*}
Запишем их разность, учитывая
\begin{equation*}
u'''(x+h/2) - u'''(x-h/2) = h u^{(4)}(x) + O(h^2).
\end{equation*}
Получим
\begin{equation*}
u'(x+h/2) - u'(x-h/2) = \frac{u(x+h) - 2 u(x) + u(x-h)}{h} - \frac{h^3}{24} u^{(4)}(x) + O(h^4).
\end{equation*}
Подставим последнее выражение в~\cref{eq:fdm_laplace_init}:
\begin{equation}
\label{eq:fdm_laplace_init}
u''(x) = -\frac{-u(x+h) + 2 u(x) - u(x-h)}{h^2} + O(h^2).
\end{equation}
Перенесём это выражение на равномерную сетку и окончательно запишем разностую схему второго порядка
для уравнения \cref{eq:poisson1d}
\begin{equation}
    \label{eq:poisson1d_fdm}
    \frac{-u_{i-1} + 2u_{i} - u_{i+1}}{h^2} = f_i, \qquad i=\overline{1,N-1}.
\end{equation}
Здесь $\{f_i\}$ -- известный сеточный вектор, определяемый через известную
аналитическую функцию $f(x)$ в правой части уравнения \cref{eq:poisson1d} как
\begin{equation}
    \label{eq:poisson1d_fdm2}
    f_i = f(x_i).
\end{equation}

\subsubsection{Сборка СЛАУ}
Линейные уравнения \cref{eq:poisson1d_fdm} относительно неизвестных $u_i$ составляют систему вида
\begin{equation*}
    \sum_{j=0}^{N} A_{ij}\,u_j = b_i, \qquad i=\overline{1,N-1}
\end{equation*}
с матричными коэффициентами
\begin{equation}
    \label{eq:poisson1d_fdm_lhs}
    A_{ij} = \begin{cases}
        2/h^2,  &\quad i=\overline{1,N-1}, \, j=i;\\
        -1/h^2, &\quad i=\overline{1,N-1}, \, j=i-1;\\
        -1/h^2, &\quad i=\overline{1,N-1}, \, j=i+1;\\
        0,      &\quad \text{иначе}.
    \end{cases}
\end{equation}
и правой частью
\begin{equation}
    \label{eq:poisson1d_fdm_rhs}
    b_i = f_i,   \quad i=\overline{1,N-1}.
\end{equation}
Всего в этой системе $N-1$ уравнение для $N+1$ неизвестной.
Для получения однозначного ответа не хватает двух уравнений, которые
должны быть сформулированы из граничных условий (для правой и левой границ).

\subsubsection{Учёт граничных условий}
\subsubsubsection{Граничные условия первого рода}

Рассмотрим условия первого рода вида:
\begin{equation}
	\label{eq:poisson1d_bc}
	\begin{cases}
        u(a)=u_a,\\[5pt]
        u(b)=u_b.\\
	\end{cases}
\end{equation}
Эти условия для сеточного вектора приобретут вид:
\begin{equation*}
u_0 = u_a, \qquad u_{N} = u_b.
\end{equation*}
В терминах матричных коэффициентов~\cref{eq:poisson1d_fdm_lhs,eq:poisson1d_fdm_rhs}
эти условия можно переписать в виде:
\begin{equation}
\begin{array}{ll}
A_{00} = 1, & b_{0} = u_a; \\
A_{NN} = 1, & b_{N} = u_b; \\
\end{array}
\end{equation}

\subsubsubsection{Граничные условия второго рода}
Рассмотрим условия Неймана
\begin{equation}
	\label{eq:poisson1d_bc2}
	\begin{cases}
        -\left.\ddfr{u}{n}\right|_a=q_a,\\[10pt]
        -\left.\ddfr{u}{n}\right|_b=q_b,
	\end{cases}
\end{equation}
Для их аппроксимации можно бы было воспользоваться направленными
разностями \cref{eq:forwdiff}. Учтём что нормаль $\vec n$ направлена наружу (против оси $x$ на левой границе и по оси $x$ -- на правой).
Тогда можно записать
\begin{equation}
- \frac{u_0 - u_1}{h} = q_a, \quad
- \frac{u_N - u_{N-1}}{h} = q_b,
\end{equation}
В терминах матричных коэффициентов получим
\begin{equation}
\label{eq:poisson1d_bc2_o1}
\begin{array}{lll}
A_{00} = 1/h & A_{01} = -1/h & b_0 = -q_a \\
A_{NN} = 1/h & A_{N,N-1} = -1/h & b_N = -q_b
\end{array}
\end{equation}

Следует однако заметить, что формула \cref{eq:forwdiff} имеет первый порядок
точности. Поэтому общий порядок расчётной схемы также будет первый.
Чтобы аппроксимировать граничные условия второго рода вторым порядком
воспользуемся уравнением \cref{eq:poisson1d_fdm} для
аппроксимации оператора Лапласа в граничной (левой) точке:
\begin{equation*}
u''_0 = \frac{u_{-1} - 2 u_{0} + u_{1}}{h^2}.
\end{equation*}
Здесь $u_{-1}$ -- это фиктивный узел, находящийся вне области расчёта.
Для нахождения его значения
аппроксимируем граничное условиe \cref{eq:poisson1d_bc2}
с использованием симметричной схемы \cref{eq:symdiff}:
\begin{equation*}
-\frac{u_{-1} - u_{1}}{2h} = q_a \hence
u_{-1} = u_{1} - 2 h q_a.
\end{equation*}
Используя два последних соотношения
запишем аппроксимированное уравнение Пуассона в граничной точке:
\begin{equation}
\label{eq:poisson1d_bc2left_o2}
-u''_0 = f_0 \hence
\frac{u_{0} - u_{1}}{h^2} = -\frac{q_a}{h} + \frac{1}{2} f_0.
\end{equation}
Переписывая в терминах матричных коэффициентов получим (сохраняя $h^2$ в знаменателе, чтобы
выражение имело тот же порядок, что и выражение для внутренних узлов \cref{eq:poisson1d_fdm}):
\begin{equation}
\label{eq:poisson1d_bc2_o2}
\begin{array}{lll}
A_{00} = 1/h^2 & A_{01} = -1/h^2 & b_0 = -\sfrac{q_a}{h} + \sfrac{f_0}{2} \\
A_{NN} = 1/h^2 & A_{N,N-1} = -1/h^2 & b_N = -\sfrac{q_b}{h} + \sfrac{f_N}{2}
\end{array}
\end{equation}
Ещё раз заметим, что все аппроксимационные формулы,
использованные для получения \cref{eq:poisson1d_bc2_o2}
имели второй порядок точности, поэтому
и само это выражение имеет второй порядок (в отличии от альтернативного \cref{eq:poisson1d_bc2_o1}).


\subsubsubsection{Граничные условия третьего рода}
Аппроксимацию условий третьего рода
\begin{equation}
	\label{eq:poisson1d_bc2}
	\begin{cases}
        -\left.\ddfr{u}{n}\right|_a=\alpha_a u(a) + \beta_a,\\[10pt]
        -\left.\ddfr{u}{n}\right|_b=\alpha_b u(b) + \beta_b,
	\end{cases}
\end{equation}
можно получить по аналогии с аппроксимацией
условий второго рода, заменив $q$ на $\alpha u + \beta$.
В частности, для левого узла
аппроксимация второго порядка по аналогии с \cref{eq:poisson1d_bc2left_o2}
примет вид
\begin{equation*}
\label{eq:poisson1d_bc3left_o2}
-u''_0 = \alpha_a u_0 + \beta \hence
\frac{(1 + \alpha_a h) u_{0} - u_{1}}{h^2} = -\frac{\beta}{h} + \frac{1}{2} f_0.
\end{equation*}
А на матричном уровне получим
\begin{equation}
\label{eq:poisson1d_bc3_o2}
\begin{array}{lll}
A_{00} = (1 + \alpha_a h)/h^2 & A_{01} = -1/h^2 & b_0 = -\sfrac{\beta_a}{h} + \sfrac{f_0}{2} \\
A_{NN} = (1 + \alpha_b h)/h^2 & A_{N,N-1} = -1/h^2 & b_N = -\sfrac{\beta_b}{h} + \sfrac{f_N}{2}
\end{array}
\end{equation}

\subsubsubsection{Периодические граничные условия}

Для удовлетворения периодических условий
можно на сеточном уровне 
закольцевать сетку, переименовав
узел с индексом ${N}$ в узел $0$, узел ${N+1}$ в $1$  и т.п.
В сеточном векторе тогда останется $N$ неизвестных.
Уравнения для крайних точек примут вид
\begin{equation}
    \label{eq:poisson1d_per0}
    \frac{-u_{N-1} + 2u_{0} - u_{1}}{h^2} = f_0, \qquad 
    \frac{-u_{N-2} + 2u_{N-1} - u_{0}}{h^2} = f_{N-1}.
\end{equation}

\subsubsection{Двумерная задача}
Двумерная задача Пуассона в частных производных будет иметь вид
\begin{equation}
    \label{eq:poisson_2d}
   -\left(\dfrq{u}{x} + \dfrq{u}{y}\right) = f(x, y).
\end{equation}

Разностную схему для такого уравнения удобно писать и использованием парного индекса $(i,j)$, где $i$ -- индекс вдоль оси $x$, $j$ -- индекс вдоль оси $y$.
При этом сеточные вектора должны по прежнему оставаться линейными массивами.
Функцию, переводящую парный индекс в сквозной (линейный), обозначим как
\begin{equation}
    \label{eq:poisson_fdm_2d_ij2k}
    k[i, j] = i + (n_x+1) j
\end{equation}
$n_x$ -- количество ячеек сетки в направлении x.
Обратные функции, переводящие сквозной индекс в пару $i,j$, имеют вид
\begin{equation}
    \label{eq:poisson_fdm_2d_k2ij}
    \begin{array}{ll}
        i[k] = {\rm mod}\left(k, \left(n_x+1\right)\right), & \text{// остаток от деления},\\
        j[k] = \lfloor k / \left(n_x+1\right)\rfloor,       &\text{// целая часть от деления}.\\
    \end{array}
\end{equation}

С учётом введённых обозначений разностная схема для уравнения \cref{eq:poisson_2d} примет вид
\begin{equation}
   \label{eq:poisson_fdm_2d}
   \frac{-u_{k[i-1,j]} + 2 u_{k[i, j]} - u_{k[i+1,j]}}{h_x^2} +
   \frac{-u_{k[i,j-1]} + 2 u_{k[i, j]} - u_{k[i,j+1]}}{h_y^2} =
   f_{k[i,j]},
\end{equation}

\subsubsection{Верификация порядка аппроксимации}
\label{sec:compute-appr}

Порядок аппрокцимации показывает скорость
приближения численного решения к точному с уменьшением сетки.
Поэтому для подтверждения порядка необходимо
\begin{itemize}
\item Знать точное решение,
\item Уметь вычислять функционал (норму, $||\cdot||$), характеризующий отклонение точного решения от численного,
\item Сделать несколько расчётов на сетках с разной $N$  и заполнить таблицу $||\{u_i - u^e(x_i)\}||(N)$,
\item На основе этой таблицы построить график в логарифмических осях и по углу наклона кривой сделать вывод о порядке аппроксимации.
\end{itemize}

Выберем произвольную функцию $u^e$ (достаточно сильно изменяющуюся на целевом отрезке $[a,b]$).

Далее путём прямого вычисления определим параметры задачи $f$, $u_a$, $u_b$ такие,
для которых функция $u^e$ является точным решением задачи \cref{eq:poisson1d}, \cref{eq:poisson1d_bc}.

Зададимся числом разбиений $N$ и решим задачу для выбранным параметров.
В результате определим сеточный вектор численного решения $\{u_i\}$.

В качестве нормы выберем стандартное отклонение. В интегральном виде для многомерной функции $y(\vec x)$
в области $\vec x\in D$ оно имеет вид
\begin{equation}
    \label{eq:norm2_common}
    ||y(\vec x)||_2 = \sqrt{\frac{1}{|D|}\int_{D} y(\vec x)^2 \, d\vec x}.
\end{equation}
Упрощая до одномерного случая
\begin{equation*}
    ||y(x)||_2 = \sqrt{\frac{1}{b-a}\int_{a}^{b} y(x)^2 \, dx}.
\end{equation*}

Вычислим этот интеграл численно на введённой ранее равномерной сетке $\{x_i\}$:
\begin{equation*}
    ||\{y_i\}||_2 = \sqrt{\frac{1}{b-a}\sum_{i=0}^{N} w_i y_i^2},
\end{equation*}
где $\{w_i\}$ -- вес (или "площадь влияния") $i$-ого узла:
\begin{equation*}
    w_i = \begin{cases}
        h/2, &\quad i=0, N;\\
        h, &\quad i=\overline{1,N-1},
    \end{cases}
\end{equation*}
такая что
\begin{equation*}
    \sum_{i=0}^{N} w_i = b-a.
\end{equation*}

Окончательно среднеквадратичная норма отклонения численного решения от точного запишется в виде
\begin{equation}
    \label{eq:poisson1d_fdm_norm}
    ||\{u_i - u^e(x_i)\}||_2 = \sqrt{\frac{1}{b-a}\sum_{i=0}^{N} w_i \left(u_i - u^e_i\right)^2}.
\end{equation}

Пример программы, реализующей этот алгоритм смотри в \secref{sec:poisson1d_prog}.
